\chapter{术语表}
\chapterlead{按主题分组的术语，同组内按出现顺序排列。英文原词放在括号里，检索文档与报错时用英文更有效。}

\section{文件与磁盘}

\begin{booktable}{\textwidth}{J{0.5} J{1.5}}
\toprule
术语 & 说明 \\
\midrule
\endhead
分区（partition） & 磁盘上的独立区域，各自容纳一个文件系统。第\ref{ch:install}章 \\
挂载（mount） & 把文件系统接到目录树上某个位置。第\ref{ch:install}章 \\
校验和（checksum） & 用于验证数据完整性的指纹，例如 SHA-256。第\ref{ch:install}章 \\
路径（path） & 文件在目录树中的位置；以斜线开头为绝对路径。第\ref{ch:files-and-dirs}章 \\
符号链接（symlink） & 内容为另一条路径的文件，可跨分区。第\ref{ch:files-and-dirs}章 \\
软链接 & 符号链接的别称。第\ref{ch:files-and-dirs}章 \\
索引节点（inode） & 文件系统记录文件元数据与数据块位置的结构。第\ref{ch:files-and-dirs}章 \\
虚拟文件系统（procfs） & 由内核生成、反映运行时状态的目录树，挂在 \texttt{/proc}。第\ref{ch:unix-culture}章 \\
\bottomrule
\end{booktable}

\section{命令行}

\begin{booktable}{\textwidth}{J{0.5} J{1.5}}
\toprule
术语 & 说明 \\
\midrule
\endhead
命令行界面（CLI） & 以文本输入输出为主的交互方式。第\ref{ch:first-terminal}章 \\
伪终端（pty） & 终端模拟器与 shell 之间的通道。第\ref{ch:first-terminal}章 \\
环境变量（environment variable） & 影响程序行为的键值，例如 \texttt{PATH}。第\ref{ch:first-terminal}章 \\
管道（pipe） & 把左侧程序的标准输出接到右侧程序的标准输入。第\ref{ch:unix-culture}章 \\
标准输出（stdout） & 程序的正常输出通道，编号 1。第\ref{ch:unix-culture}章 \\
标准错误（stderr） & 程序的错误输出通道，编号 2。第\ref{ch:unix-culture}章 \\
\bottomrule
\end{booktable}

\section{权限、进程与服务}

\begin{booktable}{\textwidth}{J{0.5} J{1.5}}
\toprule
术语 & 说明 \\
\midrule
\endhead
权限位（mode） & 文件或目录的读、写、执行许可，九位三组。第\ref{ch:permissions}章 \\
系统调用（syscall） & 程序请求内核服务的入口。第\ref{ch:what-a-computer-does}章 \\
内核（kernel） & 操作系统中直接访问硬件、管理进程与内存的核心部分。第\ref{ch:what-a-computer-does}章 \\
信号（signal） & 内核发给进程的简短消息，如 \texttt{SIGTERM}。第\ref{ch:processes}章 \\
守护进程（daemon） & 长期在后台运行、无终端的进程。第\ref{ch:processes}章 \\
目标（target） & systemd 中一组单元的集合，例如 \texttt{multi-user.target}。第\ref{ch:processes}章 \\
用户单元（user unit） & 属于某个登录用户的 systemd 单元。第\ref{ch:processes}章 \\
轮转（rotation） & 日志按大小或时间切分并清理旧文件的机制。第\ref{ch:monitoring}章 \\
容器（container） & 共享宿主内核的隔离进程。第\ref{ch:containers}章 \\
\bottomrule
\end{booktable}

\section{NixOS}

\begin{booktable}{\textwidth}{J{0.5} J{1.5}}
\toprule
术语 & 说明 \\
\midrule
\endhead
命令式（imperative） & 由操作序列改变状态，最终形态取决于操作历史。第\ref{ch:why-nixos}章 \\
声明式（declarative） & 由描述算出目标状态，结果与操作的先后无关。第\ref{ch:why-nixos}章 \\
发行版（distribution） & 内核、工具、包管理器、默认配置与支持渠道的打包。第\ref{ch:what-a-computer-does}章 \\
可复现（reproducible） & 同样的输入得到同样的产物。第\ref{ch:nixos-model}章 \\
散列（hash） & 由内容算出的定长指纹，用于命名 store 路径。第\ref{ch:nixos-model}章 \\
导出（derivation） & Nix 的构建配方：输入、脚本、输出路径。第\ref{ch:nixos-model}章 \\
闭包（closure） & 一个程序运行所需的全部 store 路径。第\ref{ch:nixos-model}章 \\
代（generation） & 一次系统重建留下的完整状态，可整体切换与回滚。第\ref{ch:nixos-model}章 \\
原子性（atomic） & 操作要么全部生效，要么完全不生效。第\ref{ch:upgrade-rollback}章 \\
订阅通道（channel） & Nix 传统升级方式，指向 nixpkgs 的某个分支。第\ref{ch:upgrade-rollback}章 \\
\bottomrule
\end{booktable}
